\documentclass[aspectratio=169,10pt]{beamer}

% Academic Monochrome / Plain White Theme
\usetheme{default}
\usecolortheme{dove}

\setbeamertemplate{navigation symbols}{}
\setbeamertemplate{footline}{
    \hbox{%
    \begin{beamercolorbox}[wd=.8\paperwidth,ht=2.5ex,dp=1ex,left]{author in head/foot}%
        \hspace*{1.5em}\scriptsize SQLens: Educational SQL Compiler $|$ BCSE307P Review 2
    \end{beamercolorbox}%
    \begin{beamercolorbox}[wd=.2\paperwidth,ht=2.5ex,dp=1ex,right]{date in head/foot}%
        \scriptsize \insertframenumber{} / \inserttotalframenumber\hspace*{1.5em}
    \end{beamercolorbox}}%
    \vskip0pt%
}

\usepackage{booktabs}
\usepackage{tabularx}
\usepackage{tikz}
\usetikzlibrary{shapes,arrows.meta,positioning}

\title{\textbf{SQLens: An Educational SQL Compiler and Query Analysis Visualizer}}
\subtitle{Compiler Design Laboratory (BCSE307P) --- Review 2: Implementation Progress Review}
\author{\textbf{Mishti Parikshit Gandhi} (Reg. No. 24BCE0881)\\Laboratory Slot: L49 + L50\\[1ex]Faculty Reviewer: \textbf{[Faculty Name / Reviewer]}}
\institute{School of Computer Science and Engineering (SCOPE)\\Vellore Institute of Technology, Vellore}
\date{October 2026}

\begin{document}

% Slide 1: Title
\begin{frame}[plain]
    \titlepage
    \note{Good morning respected faculty evaluator. I am Mishti Parikshit Gandhi, registration number 24BCE0881, presenting Review 2 for my Compiler Design Laboratory project titled SQLens.}
\end{frame}

% Slide 2: Project Overview
\begin{frame}{Project Overview: Concept, Purpose \& Current Status}
    \begin{columns}[T]
        \begin{column}{0.48\textwidth}
            \textbf{What is SQLens \& What Does It Do?}
            \begin{itemize}\small
                \item \textbf{Domain-Specific Compiler:} Educational SQL compiler built in C++ (Flex, Bison) for a declarative MiniSQL language.
                \item \textbf{Bridging Core Fields:} Connects Compiler Design (lexing, LALR parsing, ASTs, type checks) with DBMS (relational algebra, optimization, CSV scans).
                \item \textbf{End-to-End Transparency:} Opens the internal pipeline into 10 observable compilation stages.
                \item \textbf{Zero Mock Data:} All outputs and trees are generated directly by \texttt{minisql.exe}.
            \end{itemize}
        \end{column}
        \begin{column}{0.48\textwidth}
            \textbf{Review 2 Implementation Status}
            \begin{itemize}\small
                \item \textbf{Phase 1 Frontend (100\%):} Lexer (Flex), Parser (Bison), AST, Symbol Table, and Semantic Analyzer verified.
                \item \textbf{Phase 2 Backend (100\%):} Relational IR, Rule-based Optimizer, and in-memory CSV Execution Engine verified.
                \item \textbf{Automated Test Suite (100\%):} 26/26 automated unit and integration tests passing.
                \item \textbf{Web Visualizer (~80\%):} React 18 + Node.js API bridge with animated stepper and stage inspector.
            \end{itemize}
        \end{column}
    \end{columns}
    \note{SQLens exposes all 10 stages of query compilation. Traditional compilers are treated as black boxes. For Review 2, our entire C++ frontend and backend are 100\% complete and tested.}
\end{frame}

% Slide 3: Revised Problem Statement
\begin{frame}{Revised Problem Statement \& Pedagogical Need}
    \textbf{The Pedagogical Problem in Traditional Instruction:}
    \begin{itemize}\small
        \item Compiler courses focus on imperative languages (C, Java), leaving database query engines as opaque black boxes.
        \item DBMS courses teach relational algebra and physical plans theoretically without showing compiler frontend transformations.
        \item Monolithic industrial toolchains (GCC, LLVM) prevent students from inspecting intermediate representations and stage diagnostics.
    \end{itemize}
    \vspace{0.8em}
    \textbf{Formal Problem Statement for Review 2:}
    \begin{quote}\small\itshape
        “Compiler Design concepts can be difficult for students to understand because the internal transformations performed during compilation are not always visible. SQLens addresses this problem by implementing a lightweight MiniSQL compiler and exposing its query-processing pipeline through an interactive visual interface. It enables learners to inspect tokens, syntax structures, semantic analysis, intermediate representation, optimization, execution, and errors through one integrated workflow.”
    \end{quote}
    \note{Our revised problem statement addresses the disconnect between theoretical compiler design and visible database execution.}
\end{frame}

% Slide 4: Objectives
\begin{frame}{Project Objectives \& Scope of Implementation}
    \begin{enumerate}\small
        \item \textbf{Lexical \& Syntax Parsing:} Implement a formal MiniSQL scanner with Flex and an LALR(1) context-free grammar with Bison to parse declarative queries and establish operator precedence.
        \item \textbf{AST \& Symbol Catalog:} Construct a strongly-typed Abstract Syntax Tree and maintain an in-memory catalog schema verifying relation names, attribute existence, and column types.
        \item \textbf{Semantic Validation:} Implement strict semantic analysis enforcing table resolution, column scoping, operator type-compatibility, and boolean predicate validation.
        \item \textbf{Relational Intermediate Code:} Translate validated ASTs into canonical Relational Algebra Intermediate Representation (IR) trees comprising Scan, Filter, Project, Sort, and Limit operators.
        \item \textbf{Rule-Based Query Optimization:} Implement algebraic query transformations including Filter Pushdown, Late Projection (projection reordering), and Limit Early-Stopping.
        \item \textbf{CSV Execution \& Fail-Stop Errors:} Execute physical relational plans against CSV datasets and enforce fail-stop error isolation to halt execution at the exact stage of failure.
    \end{enumerate}
    \note{These 6 measurable objectives are directly implemented in our C++ codebase and validated by our test harness.}
\end{frame}

% Slide 5: Scope and Limitations
\begin{frame}{Scope \& Limitations: Supported Syntax vs Out-of-Scope}
    \begin{columns}[T]
        \begin{column}{0.48\textwidth}
            \textbf{Supported MiniSQL Language (Verified in Code)}
            \begin{itemize}\footnotesize
                \item \textbf{Statements:} Declarative \texttt{SELECT ... FROM ... [WHERE ...] [ORDER BY ...] [LIMIT ...]};
                \item \textbf{Relations:} Single catalog table \texttt{students} (\texttt{id}, \texttt{name}, \texttt{cgpa}, \texttt{age}).
                \item \textbf{Projections:} Single/multiple comma-separated columns.
                \item \textbf{Predicates:} Comparisons (\texttt{=}, \texttt{!=}, \texttt{<}, \texttt{>}, \texttt{<=}, \texttt{>=}), logical \texttt{AND}, \texttt{OR}, and parenthesized expressions.
                \item \textbf{Ordering \& Limits:} \texttt{ORDER BY} [ASC$|$DESC] (including unprojected sort keys), \texttt{LIMIT <int>}.
                \item \textbf{Storage:} Sequential streaming CSV parsing.
            \end{itemize}
        \end{column}
        \begin{column}{0.48\textwidth}
            \textbf{Explicitly Out-of-Scope / Phase 3 Planned}
            \begin{itemize}\footnotesize
                \item \textbf{Multi-Table Operations:} Joins (\texttt{INNER JOIN}, \texttt{CROSS JOIN}) reserved for Phase 3.
                \item \textbf{Aggregations:} \texttt{COUNT()}, \texttt{SUM()}, \texttt{AVG()}, \texttt{GROUP BY}, and \texttt{HAVING} are excluded.
                \item \textbf{DML / DDL:} \texttt{INSERT}, \texttt{UPDATE}, \texttt{DELETE}, \texttt{CREATE TABLE} excluded (focus is strictly on query compilation).
                \item \textbf{Nested Queries:} Subqueries in \texttt{WHERE} or \texttt{FROM}.
                \item \textbf{Dynamic Schemas:} Arbitrary user-uploaded schemas.
            \end{itemize}
        \end{column}
    \end{columns}
    \note{We maintain academic integrity by clearly defining our boundaries: MiniSQL focuses deeply on complete compiler transparency over single relations.}
\end{frame}

% Slide 6: Overall System Architecture
\begin{frame}{Overall System Architecture \& Component Flow}
    \begin{columns}[T]
        \begin{column}{0.32\textwidth}
            \textbf{1. Web Interface Layer}
            \begin{itemize}\footnotesize
                \item React 18 + Vite (Port 5173).
                \item Express API Bridge (Port 3001).
                \item ChildProcess CLI execution of \texttt{minisql.exe --json}.
                \item Pure structured JSON response.
            \end{itemize}
        \end{column}
        \begin{column}{0.35\textwidth}
            \textbf{2. 10-Stage C++ Compiler}
            \begin{itemize}\footnotesize
                \item [01] Source Query Validation
                \item [02] Flex Lexical Scanner
                \item [03] Bison LALR(1) Parser
                \item [04] Typed AST Construction
                \item [05] Symbol Table Catalog
                \item [06] Semantic Type Checker
                \item [07] Relational IR Operator Tree
                \item [08] Rule-Based Optimizer
                \item [09] CSV Streaming Executor
                \item [10] Query Results \& Traces
            \end{itemize}
        \end{column}
        \begin{column}{0.30\textwidth}
            \textbf{3. Storage \& Diagnostics}
            \begin{itemize}\footnotesize
                \item Data: \texttt{data/students.csv}.
                \item In-memory typed evaluation.
                \item Strict fail-stop error isolation.
                \item Nanosecond stage profiling.
            \end{itemize}
        \end{column}
    \end{columns}
    \note{The system architecture has 3 clean tiers: the React/Node visualizer, our 10-stage native C++ compiler pipeline, and physical CSV storage.}
\end{frame}

% Slide 7: Lexical and Syntax Analysis
\begin{frame}{Frontend Compilation: Lexical Scanning \& Syntax Parsing}
    \begin{columns}[T]
        \begin{column}{0.48\textwidth}
            \textbf{Lexical Analysis: Flex (\texttt{lexer.l})}
            \begin{itemize}\footnotesize
                \item Scans character stream using DFA regular expressions.
                \item Emits \texttt{TokenRecord}: token name, original lexeme, token classification, line number.
                \item Catches illegal characters (e.g. \texttt{@}, \texttt{\$}) via fallback rules, triggering immediate fail-stop termination.
            \end{itemize}
        \end{column}
        \begin{column}{0.48\textwidth}
            \textbf{Syntax Analysis: Bison (\texttt{parser.y})}
            \begin{itemize}\footnotesize
                \item LALR(1) context-free grammar with explicit operator precedence (\texttt{\%left OR}, \texttt{\%left AND}, \texttt{\%left comp\_ops}).
                \item Semantic actions build AST nodes bottom-up.
                \item Enabled \texttt{\%define parse.error detailed} for descriptive syntax error reports.
            \end{itemize}
        \end{column}
    \end{columns}
    \vspace{0.8em}
    \textbf{Real Token Stream Example:} \texttt{SELECT name, cgpa FROM students WHERE cgpa > 8.0;}
    \begin{block}{}\footnotesize
        \texttt{[SELECT, KEYWORD] $\to$ [IDENTIFIER(name)] $\to$ [COMMA] $\to$ [IDENTIFIER(cgpa)] $\to$ [FROM, KEYWORD] $\to$ [IDENTIFIER(students)] $\to$ [WHERE, KEYWORD] $\to$ [IDENTIFIER(cgpa)] $\to$ [GREATER, >] $\to$ [FLOAT\_LITERAL(8.0)] $\to$ [SEMICOLON]}
    \end{block}
    \note{Flex tokenizes input into structured records, and Bison enforces operator precedence and constructs AST nodes during reduction.}
\end{frame}

% Slide 8: AST and Semantic Analysis
\begin{frame}{Frontend Compilation: AST Hierarchy \& Semantic Validation}
    \begin{columns}[T]
        \begin{column}{0.48\textwidth}
            \textbf{Abstract Syntax Tree (\texttt{ast.h}, \texttt{ast.cpp})}
            \begin{itemize}\footnotesize
                \item Base \texttt{ASTNode} with derived \texttt{SelectQueryNode} and \texttt{ExprNode} hierarchy.
                \item Decouples semantic query structure from parsing grammar and surface punctuation.
                \item Implements recursive ASCII tree rendering and JSON serialization.
            \end{itemize}
        \end{column}
        \begin{column}{0.48\textwidth}
            \textbf{Symbol Table \& Semantic Analysis}
            \begin{itemize}\footnotesize
                \item Catalog schema for \texttt{students}: \texttt{id} (INT), \texttt{name} (STRING), \texttt{cgpa} (FLOAT), \texttt{age} (INT).
                \item 6-Point Verification: (1) Table existence, (2) SELECT columns, (3) WHERE columns, (4) ORDER BY columns, (5) Operator type compatibility, (6) Boolean predicate evaluation.
            \end{itemize}
        \end{column}
    \end{columns}
    \vspace{0.6em}
    \textbf{Semantic Verification Examples (Verified in Code):}
    \begin{itemize}\footnotesize
        \item \textbf{Valid Query:} \texttt{SELECT name, cgpa FROM students WHERE cgpa > 8.0;} $\to$ \textbf{PASS}
        \item \textbf{Unknown Column:} \texttt{SELECT salary FROM students;} $\to$ \textbf{FAIL (Stage 06)}: Column 'salary' does not exist.
        \item \textbf{Type Mismatch:} \texttt{SELECT name FROM students WHERE cgpa > 'Alice';} $\to$ \textbf{FAIL (Stage 06)}: Incompatible operand types.
    \end{itemize}
    \note{The semantic analyzer verifies attribute existence and enforces strict type safety before intermediate code generation.}
\end{frame}

% Slide 9: Intermediate Representation
\begin{frame}{Intermediate Representation: Relational Algebra Trees}
    \begin{columns}[T]
        \begin{column}{0.5\textwidth}
            \textbf{Why Relational Algebra IR? (\texttt{ir.h}, \texttt{ir.cpp})}
            \begin{itemize}\footnotesize
                \item AST models surface syntax; Relational IR models relational algebra operators independent of physical execution.
                \item Canonical plan is constructed systematically from the validated AST.
                \item Provides the formal foundation for algebraic query rewrites.
            \end{itemize}
            \textbf{Implemented Operator Nodes:}
            \begin{itemize}\footnotesize
                \item \textbf{ScanNode:} Streams base tuples from storage.
                \item \textbf{FilterNode:} Evaluates predicate ($\sigma_p$) per tuple.
                \item \textbf{SortNode:} Orders tuples by key column ($\tau$).
                \item \textbf{ProjectNode:} Restricts tuple width ($\pi_A$).
                \item \textbf{LimitNode:} Enforces maximum tuple cardinality.
            \end{itemize}
        \end{column}
        \begin{column}{0.46\textwidth}
            \textbf{Canonical Plan for Test Query:}
            \begin{block}{}\footnotesize\ttfamily
                Query: SELECT name, cgpa FROM students WHERE cgpa > 8.0 ORDER BY cgpa LIMIT 3;
            \end{block}
            \begin{center}\footnotesize
            \begin{tabular}{|l|}
                \hline
                $\blacktriangledown$ \textbf{LIMIT (3)} --- Caps row count \\
                $\blacktriangledown$ \textbf{PROJECT [name, cgpa]} --- Restricts columns \\
                $\blacktriangledown$ \textbf{SORT [cgpa ASC]} --- Orders tuples \\
                $\blacktriangledown$ \textbf{FILTER [cgpa > 8.0]} --- Predicate evaluation \\
                $\blacktriangledown$ \textbf{SCAN [students]} --- Sequential table scan \\
                \hline
            \end{tabular}
            \end{center}
        \end{column}
    \end{columns}
    \note{Unlike general compilers which emit Three-Address Code, a query compiler translates ASTs into Relational Algebra Trees.}
\end{frame}

% Slide 10: Query Optimization
\begin{frame}{Query Optimization: Rule-Based Algebraic Rewrites}
    \textbf{Implemented Algebraic Rules (\texttt{optimizer.h}, \texttt{optimizer.cpp}):}
    \begin{itemize}\footnotesize
        \item \textbf{Rule 1 --- Filter Pushdown ($\sigma_p(\pi_A(R)) \to \pi_A(\sigma_p(R))$):} Evaluates filters directly adjacent to scans. If already adjacent, optimizer flags \texttt{NOT REQUIRED}.
        \item \textbf{Rule 2 --- Late Projection (Projection Reordering):} In canonical plans, PROJECT occurs before SORT. If \texttt{ORDER BY} references a column not in \texttt{SELECT}, early projection deletes the sort key! The optimizer pushes PROJECT above SORT to preserve the sort key.
        \item \textbf{Rule 3 --- Limit Early-Stopping:} When no sort is present, pushes LIMIT into the CSV scan to halt file I/O after $N$ qualifying rows.
    \end{itemize}
    \vspace{0.5em}
    \begin{columns}[T]
        \begin{column}{0.48\textwidth}
            \textbf{Canonical Plan (Before Rewrite)}
            \begin{block}{}\footnotesize\ttfamily
                LIMIT (3)\\
                $\to$ SORT (cgpa ASC) [Needs cgpa]\\
                $\to$ PROJECT [name, age] [Strips cgpa!]\\
                $\to$ FILTER (age >= 20)\\
                $\to$ SCAN (students)
            \end{block}
        \end{column}
        \begin{column}{0.48\textwidth}
            \textbf{Optimized Plan (After Late Projection)}
            \begin{block}{}\footnotesize\ttfamily
                LIMIT (3)\\
                $\to$ PROJECT [name, age] [Pushed Above Sort!]\\
                $\to$ SORT (cgpa ASC) [cgpa Retained Safely!]\\
                $\to$ FILTER (age >= 20)\\
                $\to$ SCAN (students)
            \end{block}
        \end{column}
    \end{columns}
    \note{Late projection prevents query failure when sorting by unselected attributes, reordering projection safely above sort.}
\end{frame}

% Slide 11: Query Execution
\begin{frame}{Physical Query Execution Engine \& CSV Storage}
    \begin{columns}[T]
        \begin{column}{0.48\textwidth}
            \textbf{Execution Pipeline (\texttt{executor.cpp})}
            \begin{itemize}\footnotesize
                \item \textbf{CSV Streaming:} Reads \texttt{data/students.csv} row-by-row, validating schema headers.
                \item \textbf{Type Deserialization:} Parses text fields into runtime variants (int, float, string).
                \item \textbf{Expression Evaluator:} Dynamically evaluates boolean AST expressions.
                \item \textbf{In-Memory Quicksort:} Orders tuples according to key column and sort direction.
                \item \textbf{Projection \& Limit:} Prunes unwanted attributes and caps row output.
            \end{itemize}
        \end{column}
        \begin{column}{0.48\textwidth}
            \textbf{Trace for: \texttt{SELECT name, age FROM students WHERE age >= 20 ORDER BY cgpa LIMIT 3;}}
            \begin{block}{Physical Trace}\footnotesize
                SCAN students $\to$ 5 rows loaded\\
                FILTER age >= 20 $\to$ 3 rows pass (Alice, Bob, David)\\
                SORT by cgpa ASC $\to$ David (6.9), Bob (7.8), Alice (8.5)\\
                PROJECT [name, age] $\to$ 2 columns retained\\
                LIMIT 3 $\to$ 3 tuples emitted
            \end{block}
            \textbf{Emitted Result Set:}
            \begin{itemize}\footnotesize
                \item Row 1: \texttt{David, 22}
                \item Row 2: \texttt{Bob, 21}
                \item Row 3: \texttt{Alice, 20}
            \end{itemize}
        \end{column}
    \end{columns}
    \note{Our executor performs in-memory evaluation over CSV records, with step-by-step row logging and nanosecond profiling.}
\end{frame}

% Slide 12: Error Handling
\begin{frame}{Stage-Specific Error Handling \& Fail-Stop Isolation}
    \textbf{Core Educational Principle:} Pipeline halts at the exact originating phase; later phases are never invoked.
    \vspace{0.5em}
    \begin{table}\footnotesize
    \begin{tabularx}{\textwidth}{l p{2.2cm} p{4.2cm} X}
        \toprule
        \textbf{Error Class} & \textbf{Stage} & \textbf{Sample Input Query} & \textbf{Observed Diagnostic} \\
        \midrule
        Lexical Error & Stage 02 (Flex) & \texttt{SELECT @name FROM students;} & \texttt{[LEXICAL ERROR] Unrecognized character '@' at line 1, col 8} \\
        Syntax Error & Stage 03 (Bison) & \texttt{SELECT FROM students;} & \texttt{[SYNTAX ERROR] unexpected 'FROM', expecting identifier} \\
        Semantic Error & Stage 06 (Semantic) & \texttt{SELECT salary FROM students;} & \texttt{[SEMANTIC ERROR] Column 'salary' does not exist in table} \\
        Type Error & Stage 06 (Semantic) & \texttt{SELECT name FROM students WHERE cgpa > 'Alice';} & \texttt{[TYPE ERROR] Incompatible types for '>': FLOAT and STRING} \\
        Execution Error & Stage 09 (Executor) & Missing file or non-numeric data & \texttt{[EXECUTION ERROR] CSV file not found or parse failed} \\
        \bottomrule
    \end{tabularx}
    \end{table}
    \note{This taxonomy shows that our error isolation cleanly mirrors the theoretical separation of compiler phases.}
\end{frame}

% Slide 13: Website and Interface
\begin{frame}{Interactive Web Visualizer Workbench}
    \begin{columns}[T]
        \begin{column}{0.48\textwidth}
            \textbf{Visualizer Architecture (React 18 + Node.js)}
            \begin{itemize}\footnotesize
                \item \textbf{Frontend:} React 18, TypeScript, Tailwind CSS, Lucide icons, Vite 5 (Port 5173).
                \item \textbf{Backend API Bridge:} Node.js Express server on port 3001 executing \texttt{minisql.exe --json}.
                \item \textbf{Zero Simulated Data:} 100\% powered by the real C++ binary outputting structured JSON.
                \item \textbf{Local Offline Execution:} Operates completely on localhost without cloud dependencies.
            \end{itemize}
        \end{column}
        \begin{column}{0.48\textwidth}
            \textbf{Interactive Demonstration Features}
            \begin{itemize}\footnotesize
                \item \textbf{10-Stage Pipeline Stepper:} Play, pause, step through, or reset to observe progressive compilation.
                \item \textbf{Interactive Token Inspector:} Token type badges, raw lexemes, and line coordinates.
                \item \textbf{Collapsible AST Tree:} Visual tree of clauses, operators, and literals.
                \item \textbf{Side-by-Side IR Diff:} Canonical vs. Optimized plan comparison with rationale badges.
                \item \textbf{Live Results \& Viva Guide:} Tabular CSV output and integrated oral defense revision cards.
            \end{itemize}
        \end{column}
    \end{columns}
    \note{The web workbench acts as an interactive visual debugger, making compilation observable for viva review and classroom learning.}
\end{frame}

% Slide 14: Testing and Results
\begin{frame}{Automated Testing \& Experimental Validation}
    \textbf{PowerShell Test Suite Summary (\texttt{run\_tests.ps1}):} 26/26 Tests Passing (100\% Pass Rate).
    \vspace{0.4em}
    \begin{table}\footnotesize
    \begin{tabularx}{\textwidth}{l l p{4.6cm} p{3.2cm} c}
        \toprule
        \textbf{Category} & \textbf{IDs} & \textbf{Scenario Description} & \textbf{Expected Behavior} & \textbf{Status} \\
        \midrule
        Valid Projections & T01, T02 & Single and multi-column selection & Projected columns match & PASS \\
        Predicate Filters & T03, T04 & Float (\texttt{cgpa > 8.0}) \& Int (\texttt{age >= 20}) & Correct row filtration & PASS \\
        Logical Predicates & T05, T06 & Compound boolean \texttt{AND} / \texttt{OR} & Precedence evaluation & PASS \\
        Ordering \& Limits & T07--T09 & \texttt{ORDER BY}, \texttt{LIMIT}, combined & Sorting \& row count cap & PASS \\
        Late Projection & T10, T11 & \texttt{ORDER BY} on unselected column & Projection reordering & PASS \\
        Boundary Cases & T12--T14 & \texttt{LIMIT 0}, \texttt{LIMIT > rows}, no sort & Clean early stopping & PASS \\
        Error Isolation & T17--T22 & Lexical, Syntax, Semantic, Type errors & Fail-stop stage isolation & PASS \\
        API Bridge Output & T26 & \texttt{minisql.exe --json} bridge execution & Machine-readable JSON & PASS \\
        \bottomrule
    \end{tabularx}
    \end{table}
    \note{Every feature is backed by automated tests, verifying projections, comparisons, boolean logic, late projection, and error traps.}
\end{frame}

% Slide 15: Implementation Progress
\begin{frame}{Module-Wise Implementation Progress Matrix}
    \begin{table}\footnotesize
    \begin{tabularx}{\textwidth}{l p{3.2cm} c c X}
        \toprule
        \textbf{Module / Subsystem} & \textbf{Technology / Source} & \textbf{Phase} & \textbf{Review 2 Status} & \textbf{Verification Evidence} \\
        \midrule
        Lexical Scanner & Flex (\texttt{lexer.l}) & Phase 1 & 100\% Complete & Token streaming \& error traps \\
        Syntax Parser & Bison (\texttt{parser.y}) & Phase 1 & 100\% Complete & LALR(1) reduction test suite \\
        AST Construction & C++17 (\texttt{ast.cpp}) & Phase 1 & 100\% Complete & ASCII print \& JSON export \\
        Symbol Catalog & C++17 (\texttt{symbol\_table.cpp}) & Phase 1 & 100\% Complete & Schema metadata lookup \\
        Semantic Analyzer & C++17 (\texttt{semantic\_analyzer.cpp}) & Phase 1 & 100\% Complete & 6-point semantic checklist \\
        Relational IR & C++17 (\texttt{ir.cpp}) & Phase 2 & 100\% Complete & Canonical operator trees \\
        Query Optimizer & C++17 (\texttt{optimizer.cpp}) & Phase 2 & 100\% Complete & Filter pushdown \& Late projection \\
        CSV Executor & C++17 (\texttt{executor.cpp}) & Phase 2 & 100\% Complete & In-memory tuple execution \\
        Fail-Stop Errors & Stage-specific traps & Phase 2 & 100\% Complete & Verified via T17--T25 tests \\
        CLI JSON Bridge & \texttt{minisql.exe --json} & Phase 2/3 & 100\% Complete & Pure structured JSON output \\
        Web Visualizer UI & React 18, TypeScript, Vite & Phase 3 & ~80\% Complete & Functional on port 5173 \\
        Multi-Table Joins & Relational Join IR & Phase 3 & Pending & Reserved for Phase 3 extension \\
        \bottomrule
    \end{tabularx}
    \end{table}
    \note{All Phase 1 and Phase 2 modules are 100\% complete. Phase 3 web visualization is 80\% complete, with joins reserved for the final review.}
\end{frame}

% Slide 16: Review 2 Rubric Mapping
\begin{frame}{Review 2 Assessment Rubric \& Evidence Mapping}
    \textbf{Mapping Implementation Evidence to the Official Marking Criteria (25 Marks Total):}
    \vspace{0.4em}
    \begin{table}\footnotesize
    \begin{tabularx}{\textwidth}{l c p{6.2cm} X}
        \toprule
        \textbf{Evaluation Criteria} & \textbf{Marks} & \textbf{Implemented Project Evidence} & \textbf{Review Status} \\
        \midrule
        Implementation Progress & 7 & Complete C++ pipeline (Lexer $\to$ Parser $\to$ AST $\to$ Semantic $\to$ IR $\to$ Optimizer $\to$ Executor) and web visualizer. & Demonstrated live \\
        Technical Correctness & 5 & Verified LALR(1) grammar, AST construction, type-checking, and CSV query evaluation across 26 tests. & 26/26 Tests Passing \\
        Compiler Concept Application & 4 & Application of DFAs, CFGs, AST traversal, symbol table lookup, IR operator trees, and algebraic rewrites. & Demonstrated across 10 stages \\
        Code Quality \& Modularity & 3 & Modular C++17 header/source separation, clean OOP inheritance, proper memory cleanup, and typed React components. & Verified in repository \\
        Testing \& Test Cases & 3 & Systematic test suite (\texttt{run\_tests.ps1}) testing projections, filters, logic, ordering, limits, and errors. & 100\% Automated pass \\
        Individual Understanding & 3 & Deep grasp of compiler internals, data structures, relational algebra mappings, and oral viva readiness. & Prepared for oral viva \\
        \bottomrule
    \end{tabularx}
    \end{table}
    \note{We have mapped our deliverables directly to the 6 criteria in the official laboratory instruction manual.}
\end{frame}

% Slide 17: Related Work
\begin{frame}{Related Academic Work in Compiler Education (2020--2025)}
    \begin{table}\footnotesize
    \begin{tabularx}{\textwidth}{l l p{4.6cm} X}
        \toprule
        \textbf{Authors \& Year} & \textbf{Venue / Journal} & \textbf{Key Findings \& Contribution} & \textbf{Comparison with SQLens} \\
        \midrule
        Stamenkovi\'c et al. (2020) & \textit{CAEE} (Wiley) & Evaluated 15 compiler simulation tools; identified lack of unified end-to-end visualization. & SQLens provides unified 10-stage pipeline with web visualization. \\
        del Vado V\'irseda (2021) & \textit{ACM ITiCSE} & Proposed learning compiler design from implementation to theory using educational grammars. & SQLens applies implementation-first approach to relational SQL. \\
        del Vado V\'irseda (2022) & \textit{ACM SIGCSE} & Developed ITT interactive tutoring tool for visualizing compiler theory from implementation. & SQLens extends tutoring to database query compilers and relational IR. \\
        Hiranishi et al. (2022) & \textit{IEEJ Trans. EIS} & Built web-based visual compiler for computer education to visualize execution flow. & SQLens pairs web visualizer with a native compiled C++ core engine. \\
        Stamenkovi\'c \& Jovanovi\'c (2024) & \textit{IEEE TLT} & Showed that web-based compiler systems significantly enhance student engagement. & SQLens confirms the efficacy of browser-accessible interactive compiler tools. \\
        Koitz-Hristov et al. (2025) & \textit{ACM SIGCSE TS} & Introduced VisOpt for visualizing compiler optimization passes in CS education. & SQLens specifically focuses on relational query optimization passes. \\
        \bottomrule
    \end{tabularx}
    \end{table}
    \note{Literature from 2020 to 2025 demonstrates the pedagogical efficacy of interactive compiler visualization tools.}
\end{frame}

% Slide 18: Project-Specific Contribution
\begin{frame}{What Distinguishes SQLens? An Honest Assessment}
    \begin{columns}[T]
        \begin{column}{0.48\textwidth}
            \textbf{What SQLens Does NOT Claim}
            \begin{itemize}\small
                \item \textbf{Not a Novel Invention:} Visual compiler simulators and educational tools have been explored since the early 2000s; visualization itself is not a new research discovery.
                \item \textbf{Not a Production Database:} SQLens does not attempt to replace production engines like PostgreSQL, DuckDB, or SQLite.
                \item \textbf{Not a Speed Benchmark:} MiniSQL is designed for maximum clarity, explainability, and pedagogical inspection, not high-throughput parallel execution.
            \end{itemize}
        \end{column}
        \begin{column}{0.48\textwidth}
            \textbf{Genuine Engineering Contributions}
            \begin{itemize}\small
                \item \textbf{Cross-Disciplinary Synthesis:} Bridges Compiler Design (Flex/Bison) with DBMS query engine internals (Relational IR, Optimization, CSV Execution).
                \item \textbf{Zero-Mock Native Architecture:} The web visualizer is 100\% driven by real C++ compiled binary outputs emitted via structured JSON.
                \item \textbf{Fail-Stop Diagnostic Isolation:} Clearly reveals where queries fail (Lexical vs. Syntax vs. Semantic vs. Runtime).
                \item \textbf{Algebraic Plan Equivalence:} Side-by-side inspection of Canonical vs. Optimized relational algebra trees.
            \end{itemize}
        \end{column}
    \end{columns}
    \note{We provide an honest academic assessment: our contribution is in engineering a unified, zero-mock bridge between compilers and DBMS.}
\end{frame}

% Slide 19: Future Work & Research Potential
\begin{frame}{Future Work \& Computing Education Research Potential}
    \begin{columns}[T]
        \begin{column}{0.48\textwidth}
            \textbf{Engineering Roadmap (Phase 3 Extensions)}
            \begin{itemize}\footnotesize
                \item \textbf{Multi-Table Joins:} Implementing Nested Loop Joins and Hash Join relational operators.
                \item \textbf{Aggregations:} Adding \texttt{COUNT()}, \texttt{SUM()}, \texttt{AVG()}, and \texttt{GROUP BY} / \texttt{HAVING} clauses.
                \item \textbf{Additional Optimizer Rules:} Constant folding in WHERE predicates and dead attribute elimination.
                \item \textbf{Execution Stepping Animation:} Highlighting qualifying vs. pruned rows visually during CSV streaming.
            \end{itemize}
        \end{column}
        \begin{column}{0.48\textwidth}
            \textbf{Research Potential (ACM SIGCSE Venue)}
            \begin{itemize}\footnotesize
                \item \textbf{Target Venue:} ACM Technical Symposium on Computer Science Education (SIGCSE).
                \item \textbf{Research Question:} \textit{“Does an integrated, stage-by-stage visualization of a real SQL compiler improve undergraduate students' conceptual understanding and debugging accuracy compared to traditional instruction?”}
                \item \textbf{Prerequisite:} Rigorous empirical classroom evaluation with human student participants is required before submission.
            \end{itemize}
        \end{column}
    \end{columns}
    \note{For Phase 3 we plan multi-table joins and aggregations, and we outline a formal classroom study required for SIGCSE research.}
\end{frame}

% Slide 20: Conclusion & Demonstration Plan
\begin{frame}{Conclusion \& Live Review Demonstration Plan}
    \begin{columns}[T]
        \begin{column}{0.48\textwidth}
            \textbf{Review 2 Summary of Achievements}
            \begin{itemize}\footnotesize
                \item Fully working MiniSQL compiler built with Flex, Bison, and C++17.
                \item 10-stage compilation pipeline from source SQL to physical CSV execution.
                \item Strict fail-stop error isolation across lexical, syntax, semantic, and runtime categories.
                \item 26/26 automated tests passing with 100\% success rate.
                \item Interactive React web visualizer connected via Node.js Express bridge.
            \end{itemize}
        \end{column}
        \begin{column}{0.48\textwidth}
            \textbf{Live Review Demonstration Sequence}
            \begin{enumerate}\footnotesize
                \item \textbf{Grand Tour Query:} Run \texttt{SELECT name, age FROM students WHERE age >= 20 ORDER BY cgpa LIMIT 3;} to demonstrate all 10 stages.
                \item \textbf{Optimization Inspection:} Compare Canonical vs. Optimized IR plan and inspect Late Projection.
                \item \textbf{Semantic Fail-Stop:} Execute \texttt{SELECT salary FROM students;} (halted at Stage 06).
                \item \textbf{Lexical Error Isolation:} Execute \texttt{SELECT @name FROM students;} (halted at Stage 02).
                \item \textbf{Automated Test Suite:} Execute \texttt{.\\run\_tests.ps1} in PowerShell (26/26 passing).
            \end{enumerate}
        \end{column}
    \end{columns}
    \note{In conclusion, all Review 2 requirements are fully met. I will now proceed to the live demonstration following our 5-step sequence.}
\end{frame}

% Slide 21: References
\begin{frame}{Academic References \& Project Bibliography}
    \begin{columns}[T]
        \begin{column}{0.48\textwidth}
            \begin{itemize}\tiny
                \item \textbf{[1] Stamenkovi\'c, D., et al. (2020).} Evaluation of simulation systems suitable for teaching compiler construction courses. \textit{CAEE}, 28(5), 1160--1177. DOI: 10.1002/cae.22231
                \item \textbf{[2] del Vado V\'irseda, R. (2021).} Learning Compiler Design: From the Implementation to Theory. In \textit{Proc. ACM ITiCSE 2021}, pp. 282--288. DOI: 10.1145/3456565.3460041
                \item \textbf{[3] del Vado V\'irseda, R. (2022).} ITT: An Interactive Tutoring Tool. In \textit{Proc. ACM SIGCSE 2022}, pp. 696--702. DOI: 10.1145/3478432.3499140
                \item \textbf{[4] Hiranishi, Y., et al. (2022).} Application of Web-based Visual Compiler to Computer Education. \textit{IEEJ Trans. EIS}, 142(4), 389--396.
            \end{itemize}
        \end{column}
        \begin{column}{0.48\textwidth}
            \begin{itemize}\tiny
                \item \textbf{[5] Stamenkovi\'c, D., \& Jovanovi\'c, Z. (2024).} A Web-Based Educational System for Teaching Compilers. \textit{IEEE TLT}, 17, 819--830. DOI: 10.1109/TLT.2023.3297626
                \item \textbf{[6] Koitz-Hristov, R., et al. (2025).} VisOpt -- Visualization of Compiler Optimizations for CS Education. In \textit{Proc. ACM SIGCSE TS 2025}.
                \item \textbf{[7] VIT SCOPE (2026).} Compiler Design Laboratory Project Instruction Manual (BCSE307P), VIT Vellore.
                \item \textbf{[8] Aho, A. V., et al. (2006).} Compilers: Principles, Techniques, and Tools (2nd ed.). Addison-Wesley.
            \end{itemize}
        \end{column}
    \end{columns}
    \note{Thank you for your time. I am now open to questions from the faculty evaluator.}
\end{frame}

\end{document}
